\section{Cochain Commitments: Private Witnesses And Public Artifacts}\label{sec:public-artifacts-commitments}

This section treats cochain commitments at the boundary between private witness
data and public-facing artifacts. In this section, ``artifact'' means
public-facing artifact unless explicitly qualified. The commitment interface
itself is mixed: commitment values, payloads, hashes, roots, and verification
statements are public-facing, while openings, randomness, full cochains, and
hidden Hodge data remain witness-level unless a specified protocol path exposes
them.
These public-facing objects are the scaffolding for transparency in action and
anonymity in identity: actions are exposed through auditable public artifacts,
while identities remain inside hidden representatives and public equivalence
classes.
This section realizes the abstract public observation \((\alpha,b)\) of
\cref{def:public-observation}: the public classifier value is \(\alpha\), while
accepted cochain payloads, specified partial openings, and proof-carrying
statements supply the cochain-facing datum \(b\) without presenting the hidden
representative, Hodge--Laplacian, or harmonic kernel.

\subsection{Cochain Commitments}

Cochain commitment payloads provide the public handle on hidden cochain data. The
cochain data itself lives at three different visibility levels. Fix a hidden
representative \(x\in A_\alpha\), with induced cochain complex \(C_x^\bullet\),
coboundaries \(\delta_0,\delta_1\), and Hodge--Laplacian \(\Delta_x\).
\begin{enumerate}[label=(\roman*)]
  \item The private witness level contains the actual \(C_x^1\)-data:
    the harmonic component \(h\), the noise component \(e\), the full observed
    cochain \(b=h+e\) when it is not explicitly published, the opening data
    \((m,r)\), and the hidden restriction/Hodge data determining
    \(\Delta_x\).
  \item The public payload level contains only a serialized residue edge vector
    \[
      \bar c=\delta_0 a\pmod q,
      \qquad a=\operatorname{enc}(m,r)\in C_x^0,
    \]
    together with a modulus, an edge count, and a topology hash. This is the
    cochain commitment payload. It is public as a committed \(C^1\)-level
    artifact, but it is not the private opening, not the full hidden cochain
    with Hodge decomposition, and not a presentation of \(\Delta_x\) or
    \(\ker(\Delta_x)\).
  \item The opening/proof level may reveal local restrictions or certify
    equations involving \(\bar c\). Only the stated local values, conditions,
    or proof statements become public.
\end{enumerate}
Thus the cochain-facing observation \(b\) in \cref{def:public-observation} is
the abstract datum supplied by an accepted payload, opening, or proof-carrying
statement; the commitment payload records the public residue edge vector
\(\bar c\); and the hidden cochain witness remains private unless a protocol
path explicitly exposes a specified part of it. The canonical payload is a
mathematical object with a modulus, public residue edge values, declared edge
count, and topology hash.

\begin{definition}[Cochain commitment payload]\label{def:cochain-commitment-payload}
  A cochain commitment payload is a tuple
  \[
    (q,n,c,\tau)
  \]
  consisting of a modulus \(q\), a declared edge count \(n\), a list \(c\) of
  integer edge values, and a topology hash \(\tau\).
\end{definition}

The intended commitment map is the transaction-facing coboundary map
\[
  (m,r) \longmapsto \delta(m\Vert r) \in C^1 \pmod q,
\]
stored together with the topology hash of the supporting complex. The payload
therefore preserves the accepted modulus, the edge-vector shape, the public
topology binding, and the serialized \(C^1\)-level commitment values. It does
not expose the hidden fibration representative, private transport data, the
operative Hodge--Laplacian, or the harmonic kernel.

The primary committed mathematical object is the resulting exact
one-cochain, not a uniquely recoverable degree-zero representative. If
\(a,a'\in C^0\) are two openings, then
\[
  \delta_0a=\delta_0a'
  \quad\Longleftrightarrow\quad
  a-a'\in\ker(\delta_0).
\]
Thus openings are naturally equivalent modulo the degree-zero global-section
space \(\ker(\delta_0)\). Selecting one distinguished opening would require a
separate gauge or reference anchor; none is inferred from the raw coboundary
payload.

\begin{definition}[Commitment well-formedness]\label{def:commitment-well-formedness}
  A cochain commitment payload is well formed when its modulus is the supported
  canonical modulus \(q=12289\), its declared edge count agrees with the actual
  edge-value list length, and the declared edge count is positive.
\end{definition}

This payload predicate is deliberately only a wire-shape condition. Relative
to an ambient sheaf \(S\), the stronger mathematical acceptance condition is
the committed one-cochain coherence predicate \(c\in\mathcal G_S\) from
\cref{def:committed-one-cochain-coherence}: it additionally requires coherent
typed coboundaries, binding to \(S\), and a witness
\(\operatorname{dec}_S(c)=\delta_0a\). Consequently an accepted committed
one-cochain is a cocycle by
\cref{prop:committed-exact-cochains-are-cocycles}.

\begin{proposition}[Well-formed committed payloads]\label{prop:well-formed-committed-payloads}
  A well-formed cochain commitment payload has nonempty edge data and modulus
  greater than one.
\end{proposition}

\begin{proof}
  By well-formedness, the declared edge count equals the edge-list length and
  is positive, so the edge list is nonempty. The modulus is \(12289>1\).
\end{proof}

\begin{proposition}[Topology-compatible aggregation]\label{prop:topology-compatible-aggregation}
  Commitment aggregation is only meaningful for payloads with compatible
  modulus, declared edge count, and topology hash. This is the public
  topology-binding condition consumed by the protocol layer.
\end{proposition}

\begin{proof}
  This is the definition of aggregation compatibility: both payloads must be
  well formed, and their modulus, declared edge count, and topology hash must
  agree. Without these equalities the payload-level addition operation is not
  an accepted aggregation.
\end{proof}

The security interface for this commitment layer consists of authenticated
binding of the exact \(C^1\) payload and its topology/instance identifier,
hiding for any wrapper that claims it, homomorphic addition, homomorphic scalar
multiplication, and partial-opening soundness. The raw equality
\(\delta_0a=\delta_0a'\) supplies only the quotient relation
\(a-a'\in\ker\delta_0\); it is not computational binding of a unique
degree-zero opening. The blinded-commitment, vector-commitment, and
cohomological-MAC interfaces are separate axiom records rather than consequences
silently bundled into the payload shape.

\begin{definition}[Blinded cochain commitment boundary]\label{def:blinded-commitment-boundary}
  Fix an accepted exact payload \(c\), with decoded value
  \(\operatorname{dec}_S(c)=\delta_0a\), and write \(s\) for a separate
  harmonic blinding scalar. A typed Hodge-blinding boundary consists of the
  following data and conditions:
  \begin{enumerate}[label=(\roman*)]
    \item \(C^1\) has the orthogonal Hodge splitting
      \[
        C^1
        =
        \mathcal H^1\oplus\image(\delta_0)\oplus\image(\delta_1^\ast),
        \qquad
        \mathcal H^1\perp\image(\delta_0)\perp\image(\delta_1^\ast).
      \]
    \item For each public topology hash \(\tau\), the blinding direction
      \(H_\tau\in\mathcal H^1\) is determined by \(\tau\).
    \item The associated blinded value has the form
      \[
        c^{\mathrm{blind}}_\tau(a,s)
        =\operatorname{dec}_S(c)+sH_\tau
        =\delta_0a+sH_\tau\in C_S^1 .
      \]
    \item The exact component \(\delta_0a\) is orthogonal to \(H_\tau\).
    \item Binding and hiding for the resulting commitment scheme are separate
      cryptographic assumptions. In particular, exactness and Hodge
      orthogonality supply the geometric decomposition but do not by
      themselves prove a distributional hiding game.
  \end{enumerate}
\end{definition}

The distinction between \(a\) and \(s\) is intentional.  The former is a
degree-zero opening, defined modulo \(\ker\delta_0\), whose coboundary gives
the accepted public payload; the
latter scales the topology-derived harmonic direction.  Moreover, the
serialized exact payload and the associated Hodge-blinded value are related
objects, not definitionally the same wire object. In particular, the modular
wire representation is interpreted by \(\operatorname{dec}_S\) before the
Hodge-space addition is formed; the paper does not identify a real cochain
space with the quotient \(C^1/qC^1\).

\begin{definition}[Vector commitment boundary]\label{def:vector-commitment-boundary}
  A vector commitment boundary is well formed if and only if all of the
  following hold:
  \begin{enumerate}[label=(\roman*)]
    \item For vertex values \(a=(a_v)\), the committed edge cochain is
      \(c=\delta a\), i.e. on an oriented edge \(e=(u,v)\),
      \[
        c_e=(\delta a)_e=R_{v\to e}a_v-R_{u\to e}a_u .
      \]
    \item An opening over a revealed set \(S\) is a pair
      \[
        o_S=(a'_S,\omega_S),
      \]
      where \(a'_S=(a'_v)_{v\in S}\) is the claimed restriction and
      \(\omega_S\) is its authentication witness.
    \item Completeness: for every revealed set \(S\) of vertices, the honest
      opening
      \[
        o_S^{\mathrm{hon}}
        =\bigl(a|_S,\operatorname{Open}_S(c)\bigr)
      \]
      satisfies \(\operatorname{Ver}_\tau(c,S,o_S^{\mathrm{hon}})=1\).
    \item Soundness: if \(o_S=(a'_S,\omega_S)\) and
      \[
        \operatorname{Ver}_\tau(c,S,o_S)=1,
      \]
      then every revealed incidence satisfies the same coboundary equation
      \(c_e=R_{v\to e}a'_v-R_{u\to e}a'_u\).
    \item Position binding: the public topology hash fixes the vertex and edge
      positions with respect to which openings are interpreted.
  \end{enumerate}
\end{definition}

\begin{definition}[Cohomological MAC boundary]\label{def:cohomological-mac-boundary}
  A cohomological MAC boundary is well formed if and only if all of the
  following hold:
  \begin{enumerate}[label=(\roman*)]
    \item The key determines a Hodge--Laplacian \(\Delta\) and harmonic
      projection \(P_H:C^1\to\ker(\Delta)\).
    \item Correct tags have the form
      \[
        T_{\Delta,\nu}(m)=P_H c(m,\nu)
      \]
      for the committed cochain \(c(m,\nu)\) determined by the message and
      nonce.
    \item Verification is the Hodge-projection check
      \[
        \operatorname{Ver}_\Delta(m,\nu,T)=1
        \Longleftrightarrow
        \Delta T=0
        \quad\text{and}\quad
        P_HT=P_Hc(m,\nu).
      \]
    \item Correctness:
      \[
        \operatorname{Ver}_\Delta(m,\nu,T_{\Delta,\nu}(m))=1.
      \]
    \item Unforgeability: producing a verifying tag without the key is exactly
      the stated hidden-Hodge projection condition for this interface.
    \item Key binding: accepted tags are bound to the chosen
      Hodge--Laplacian key.
    \item Determinism: the same key, message, and nonce determine the same tag.
  \end{enumerate}
\end{definition}

Full adversarial games for hiding, unforgeability, and partial openings belong
to the subsequent protocol paper.

The commitment-facing invariants attach to the commitment payload and, after
packaging, to the public artifact layer; they are not proximity projection data.
Payload well-formedness is a shape condition: it asserts the supported modulus,
agreement between the declared edge count and the edge-value list, and a
nonempty edge list. It is not a cryptographic uniqueness theorem and it does
not assert that the cochain vector is nonzero. For example, the canonical zero
payload is well formed whenever the declared edge count is positive.

Topology compatibility is a separate public invariant. Two committed cochains
can be aggregated only when their moduli, declared edge counts, and topology
hashes agree. The weaker same-topology relation records equality of topology
hashes alone. Authentication of the exact \(C^1\) value and topology-binding
claims are supplied by the commitment-security axioms, not by the raw payload
data; those axioms do not silently upgrade the coboundary to unique
degree-zero-message binding.

At the artifact level, a hashed committed cochain is a public cochain payload
together with digest data and a binding assertion. A proof-carrying committed
cochain adds a proof system, a statement, and a verification condition. The
well-formedness predicate for the proof-carrying artifact includes
well-formedness of the hashed public artifact; well-formedness of that hashed
artifact in turn includes well-formedness of its underlying cochain payload. The
basic theorem exposed to this paper is that a well-formed proof-carrying
committed-cochain artifact has an underlying well-formed cochain commitment
payload.

When a sheaf context \(S\) is available, the refined artifact predicate
requires the underlying payload to satisfy \(c\in\mathcal G_S\), not merely
payload well-formedness. This refinement is inherited by the corresponding
proof-carrying and runtime-acceptance predicates. It separates two valid but
different claims: ordinary artifact well-formedness exposes a valid serialized
payload, whereas sheaf-coherent artifact well-formedness exposes an exact
committed one-cochain and therefore a one-cocycle.

\begin{proposition}[Proof-carrying artifact exposes payload]\label{prop:proof-carrying-artifact-exposes-payload}
  A well-formed proof-carrying committed-cochain artifact exposes a well-formed
  underlying cochain commitment payload.
\end{proposition}

\begin{proof}
  Unpack proof-carrying well-formedness to obtain well-formedness of the hashed
  public artifact, then unpack hashed-artifact well-formedness to obtain
  well-formedness of its cochain payload.
\end{proof}

\begin{proposition}[Sheaf-coherent proof artifact exposes an exact cocycle]
\label{prop:coherent-proof-artifact-exposes-cocycle}
  If a proof-carrying committed-cochain artifact is coherent with \(S\), then
  its underlying payload \(c\) lies in \(\mathcal G_S\) and satisfies
  \[
    \delta_1\operatorname{dec}_S(c)=0.
  \]
\end{proposition}

\begin{proof}
  The refined artifact predicate projects to \(c\in\mathcal G_S\). Apply
  \cref{prop:committed-exact-cochains-are-cocycles}.
\end{proof}

\subsection{Public Artifacts}

The public-facing stack has three mathematical levels. The fully public handle
level contains hashes, roots, topology hashes, nullifiers, public summaries, and
declared public parameters. The verifier-facing artifact level contains
compressed public data, cochain commitment artifacts, proof-carrying cochain
payloads, and certificate data. The private witness level contains the fibration
representative, transport/restriction data, full cochains and openings, and the
Hodge operator whose kernel selects the harmonic component. The paper's claims
are stated so that public handles and verifier-facing artifacts bind to hidden
witnesses without turning the hidden witness into public data. This paper
records the mathematical public/private separation required by the
sheaf-theoretic attacker-observation viewpoint; the associated security games
and runtime observation model are treated in the subsequent protocol paper.

These layers form the following public artifact stack:
\[
\begin{aligned}
  &\text{hidden/private object} \\
  &\quad\to \text{compressed public datum} \\
  &\quad\to \text{commitment/root/hash artifact} \\
  &\quad\to \text{proof-carrying artifact}.
\end{aligned}
\]
Verifier and oracle judgments are conditions evaluated on the terminal
public artifact, not further objects in the stack.
Each layer binds to or summarizes the preceding layer. A hash does not replace
the mathematical object, a commitment does not replace the witness, and a proof
artifact does not replace the proof statement. The purpose of the stack is to
let public verification refer to hidden sheaf/cochain data without exposing it.
In this sense, the stack separates what must be publicly accountable from what
must remain representative-private.

Hashes serve as public names, integrity bindings, and fixed-size pointers for
larger artifacts. A topology hash names the public combinatorial support; a
root commits to a batch or storage view; a transcript hash domain-separates a
proof statement. None of these hashes asserts validity by itself. Validity
comes from the condition or proof attached to the hashed object.

\subsection{Lossy Projection And Preserved Structure}

Storage roots, topology hashes, cochain roots, nullifiers, and proof objects are
public artifacts in the stack. Full sheaf, cochain, and fibration data remain
witness material unless a protocol path explicitly exposes them. The storage
mechanics are protocol-level concerns; the mathematical point here is that
stored public artifacts are compressed images of accepted witness data, not
replacements for the hidden object.

In the notation of \cref{def:public-observation}, \(\alpha\) is the public
classifier value obtained from a lossy public map. Its fiber is
\[
  A_\alpha=\pi^{-1}(\alpha),
\]
the set of hidden representatives compatible with the same public class. For
\(x\in A_\alpha\), the induced cochain complex \(C_x^\bullet\),
Hodge--Laplacian \(\Delta_x\), and harmonic kernel \(\ker(\Delta_x)\) are not
determined by \(\alpha\).

Compression creates verifier-facing public artifacts from hidden or full
witness data. The compression need not be injective; it must preserve enough
structure to support the relevant public statement.

A compressed public datum is allowed to forget hidden data. In fact, it must
forget enough data that it does not select a private representative. The
relevant question is not whether the public datum is injective, but which facts
survive the compression relation.

The preserved structure before the finite-code layer is:
\begin{enumerate}[label=(\roman*)]
  \item public homotopy-class equality for hidden representatives;
  \item cochain commitment payload well-formedness;
  \item topology binding and exact-one-cochain realization for accepted
    sheaf-coherent commitments;
  \item digest/root binding for compressed public data;
  \item proof-carrying validity of the public statement attached to an artifact.
\end{enumerate}

The lost structure is the private representative itself: full sheaf/cochain
witness data, hidden transport, the operative Hodge operator, private
randomness, and runtime lifecycle facts. Those may be bound to public artifacts
or certified by a proof-carrying statement, but they are not made public merely
by the presence of a hash, root, or commitment payload.

Byte parsing, allocation bounds, canonical encoding, concrete hash functions,
domain separation, storage inclusion proofs, proof verification, state updates,
and ephemeral-fiber lifecycle properties are TCB requirements. They are boundary
conditions, not part of the mathematical cochain commitment payload.

For HCP, the concrete sources for the abstract observation \((\alpha,b)\) are:
\(\alpha\) is the public classifier value, and \(b\) is the cochain-facing datum
supplied by an accepted cochain payload, opening, or proof-carrying statement.
For \(x\in A_\alpha\), private restriction data determines \(C_x^\bullet\),
\(\delta_0,\delta_1\), \(\Delta_x\), and the recovery target
\[
  h=\pi_{H_x}(b)\in\ker(\Delta_x),
\]
when \(b\in C_x^1\).

The artifacts in this section may bind to the hidden representative or certify
equations involving \(b\), but they do not present \(\Delta_x\),
\(\ker(\Delta_x)\), or the private fibration data. Once HCP has selected or
certified the hidden Hodge data, the fibration ID is
\[
  \operatorname{FID}(x)=\operatorname{Spec}_{\mathrm{multi}}(\Delta_x).
\]

\subsection{Plucker Coordinates}

Let \(V\) be a finite-dimensional vector space over a field \(K\), and let
\(1\le r\le \dim V\). A private rank-\(r\) frame
\[
  \mathcal F=(w_1,\ldots,w_r)
\]
spans an \(r\)-plane
\[
  W=\langle w_1,\ldots,w_r\rangle\in \operatorname{Gr}(r,V).
\]
The Plucker embedding sends this Grassmannian point to
\[
  \operatorname{Pl}(W)
  =
  [w_1\wedge\cdots\wedge w_r]
  \in \mathbb P(\bigwedge^r V).
\]
A Plucker coordinate is the public geometric datum obtained from such a private
frame: a representative of \(\operatorname{Pl}(W)\), or a digest-bound
canonical encoding of such a representative, equipped with the relevant
projective well-formedness conditions. It is required to respect projective
rescaling and the Plucker relations; changing the private frame for the same
Grassmannian point does not change the normalized Plucker coordinate.
This is the Plucker coordinate used by the proximity projection, not the HCP
public class label \(\alpha\) fixed above.

Thus the frame \(\mathcal F\) is witness material, while the Plucker coordinate is the
verifier-facing projective datum. It is not itself the witness. It can be bound
by a hash, root, or proof-carrying statement, but the actual witness remains
the hidden frame or representative.

In the proximity specialization of \cref{sec:rs-proximity-projection}, the cochain
commitment remains the public committed payload. The Plucker coordinate enters
as witness geometry for a public-binding statement
\[
  \bigl(\operatorname{Pl}(\mathcal F),c\bigr)\in\mathcal R.
\]
For the functional encoder, the projection direction is
\[
  \operatorname{Pl}(\mathcal F)
  \xrightarrow{\ E\ }
  \underbrace{c}_{\text{accepted cochain commitment}}
  \longrightarrow
  \text{plane/RS/MCA verifier data}.
\]
A proof-carrying binding relation asserts that the accepted cochain commitment
realizes the Plucker coordinate.

\subsection*{Lean Cross-Reference}

The public artifact and commitment interfaces in this section correspond to the
following Lean declarations:
\begin{leanledger}
\begin{itemize}
  \item Plucker coordinates and public binding:
    \lean{Proximity.GrassmannianPlucker},
    \lean{Proximity.GrassmannianPlucker.Concrete}, and
    \lean{Proximity.PluckerCochain.PublicBinding};
  \item payload shape and invariants:
    \lean{Axioms.CochainCommitmentData},
    \lean{CochainCommitmentData.WellFormed},
    \lean{CochainCommitmentData.AddCompatible},
    \lean{CochainCommitmentData.edge_values_nonempty_of_wellFormed}, and
    \lean{CochainCommitmentData.modulus_gt_one_of_wellFormed};
  \item sheaf-semantic commitment invariants:
    \lean{Sheaf.CellularSheaf.SheafStateSnapshot},
    \lean{Sheaf.CellularSheaf.CochainPayloadInterpretation},
    \lean{Sheaf.CellularSheaf.AcceptedCommittedOneCochain},
    \lean{Sheaf.CellularSheaf.openingEquivalent_iff_sub_mem_ker},
    \lean{Sheaf.CellularSheaf.decoded_isExactOneCochain},
    \lean{Sheaf.CellularSheaf.decoded_isOneCocycle},
    \lean{Axioms.CommittedCochainWellFormed},
    \lean{CochainCommitmentData.CoherentWith},
    \lean{Axioms.RepresentsExactOneCochain},
    \lean{Axioms.IsOneCocycle}, and
    \lean{Axioms.oneCocycle_of_committedCochainWellFormed};
  \item commitment-security interfaces:
    \lean{Axioms.CochainCommitmentAxioms},
    \lean{Axioms.HodgeBlindingContext},
    \lean{Axioms.BlindedCommittedCochain},
    \lean{Axioms.BlindedCommittedCochain.base_representsExactOneCochain},
    \lean{Axioms.BlindedCommittedCochain.base_isOneCocycle},
    \lean{Axioms.BlindedCommittedCochain.harmonicDirection_isHarmonic},
    \lean{Axioms.BlindedCommittedCochain.exactComponent_orthogonal_harmonicDirection},
    \lean{Axioms.BlindedCommittedCochain.exists_exact_harmonic_decomposition},
    \lean{Axioms.BlindedCommitmentAxioms},
    \lean{Axioms.BlindedCommitmentWellFormed},
    \lean{Axioms.VectorCommitmentAxioms},
    \lean{Axioms.VectorCommitmentWellFormed},
    \lean{Axioms.CohomologicalMACAxioms}, and
    \lean{Axioms.CohomologicalMACWellFormed};
  \item proof-carrying artifacts:
    \lean{TCB.HashedCommittedCochain},
    \lean{TCB.ProofCarryingCommittedCochain}, and
    \lean{TCB.proofCarrying_cochain_wellFormed};
  \item sheaf-coherent artifact refinements:
    \lean{TCB.HashedCommittedCochain.CoherentWith},
    \lean{TCB.ProofCarryingCommittedCochain.CoherentWith},
    \lean{TCB.proofCarrying_cochain_coherentWith}, and
    \lean{TCB.RuntimeAcceptsSheafCommittedCochain}.
\end{itemize}
\end{leanledger}
